\section*{Abstract}
\addcontentsline{toc}{section}{Abstract}

The Nakuja N4 rocket recovery system addresses three failure modes identified in previous missions: insufficient telemetry range, uncontrolled nichrome-wire burnout during parachute deployment, and absence of pre-flight electrical status feedback. This interim report presents the design, analysis, and fabrication-stage results achieved to date.

The telemetry subsystem adopts the XBee-PRO 900HP operating at 900~MHz in transparent Universal Asynchronous Receiver-Transmitter (UART) mode, with a first-order link budget confirming a +32~dB margin at 5~km using omnidirectional dipole antennas. The ignition subsystem replaces the previous direct DC dump with Pulse-Width Modulation (PWM)-controlled IRLZ44N Metal-Oxide-Semiconductor Field-Effect Transistor (MOSFET) switching at 1~kHz; bench characterisation of 26~AWG nichrome elements demonstrated sustained energisation for approximately 30~seconds at minimum duty cycle compared to burnout within 2~seconds at DC-equivalent drive, validating the controlled-heating approach. A battery monitoring and continuity-detection circuit centred on the ADS1115 16-bit Analog-to-Digital Converter (ADC) provides pre-launch voltage and igniter-state visibility. Mechanically, a stepped polyethylene terephthalate glycol (PETG) battery-carrier architecture accommodates four 18650 cells within the 95~mm avionics bay while routing dual threaded retention rods without interference.

At the close of this interim phase, all subsystem designs have been completed and verified through 2D/3D computer-aided design (CAD), EasyEDA printed circuit board (PCB) schematics, and Proteus circuit simulation. A custom flight-computer PCB is currently in fabrication, and 3D-printed structural parts are in production. Firmware modules for telemetry, sensor interfacing, PWM control, and flight-state management have been implemented on FreeRTOS. Pending work for the next phase includes PCB assembly, open-field XBee range testing, hardware-level monitoring calibration, and pyrotechnic pop tests to validate deployment force.


